\section{El mismo control para el valor canónico entero}
Escribimos $C$ para la incidencia sitio--RC y $D_{ef}=z_e{\bf1}_{e\in f}$ para las
dos pendientes; se pueden mantener las líneas vacías con pivote uno. Sea
$T=\operatorname{diag}(q_e(1-q_e)/(1-d_e q_e))$,
$T_0=\operatorname{diag}(P_e(1-P_e)/(1-2P_e))$, $S_R=C^tTC$ y $S_0=C^tT_0 C$.
Denotamos por $P_g$ la proyección ortogonal sobre la dirección nula RC.
En el cofactor completo, Schur en ese gauge da
\[
 S_d=\operatorname{diag}h + D^tTD-D^tTC(S_R+P_g)^{-1} C^tTD .
\]
En cada máscara física, el pseudodeterminante es $\det(S_R+P_g)\det S_d$;
las columnas de $C^tTD$ son ortogonales al gauge. El soporte RC es fijo y conexo
por la brecha; su covolumen coincide en todas las activaciones. La positividad real de
ambos determinantes se sigue de $T>0$, del bloque de holguras positivas y del gauge único.

Damos las cotas con la dependencia en $\rho$, útil cuando varía con $n$. Basta tomar
$\rho\ge2$ y $\rho^2(\epsilon_A+1/j)\le\eta$ para un $\eta=\eta(M_0,r)>0$ pequeño.
En la parte 2, la norma de $H(0)$ y su constante de Lipschitz quedan acotadas por
$C(\rho^2\epsilon_A+\rho M_0/j)$ con $C$ independiente de $\rho$. El argumento
para $\Pi$ sólo usa la norma de la inversa escalar y las series exponenciales.
Por tanto $\max_e\|q_e/P_e-1\|_\rho=O_{M_0,r}(\eta)$.
También $\|T_e/T_{0,e}-1\|_\rho=O_{M_0,r}(\eta)$. Se tiene
$S_0-D_0=O(1/j)$ en norma de suma por filas; aquí $D_0$ es el jacobiano RC de
la parte 2. Se deduce $\|(S_0+P_g)^{-1}\|_\infty\le C$, y
$\|S_R-S_0\|_\infty\le C\eta$. El logaritmo de determinante RC relativo se
calcula por la serie en trazas con perturbación de norma menor que $1/2$ y cuesta
$O(j)$ en el álgebra. La inversa con gauge en el Schur queda acotada.

La norma máxima de suma por filas es $O(\rho\epsilon_A)$ en $D^tTC$,
$O(\rho)$ en $C^tTD$ y $O(\rho^2\epsilon_A)$ en $D^tTD$. En efecto, una fila diagonal
tiene masa basal $\le\epsilon_A$ y por cada sitio hay dos incidencias en cada familia de
índices (dos RC en total y dos diagonales en total). La norma de $\operatorname{diag}h-I$
es $O(\rho\epsilon_A)$. Se obtiene $\|S_d-I\|_\infty\le C\eta<1/2$. Sea $\|B\|_\Sigma$
la suma de normas de todas las entradas. Como $\sum_eP_e=j$, se tiene
$\|D^tTC\|_\Sigma=O(\rho j)$, $\|D^tTD\|_\Sigma=O(\rho^2j)$ y
$\|\operatorname{diag}h-I\|_\Sigma=O(\rho j)$. Usando
$\|BC\|_\Sigma\le\|B\|_\Sigma\|C\|_\infty$ queda $\|S_d-I\|_\Sigma=O(\rho^2j)$.
La serie de trazas de su logaritmo cuesta $O(\rho^2j)$ aun con muchas líneas.
Estas series coinciden en cada evaluación con los logaritmos reales, por continuación
desde la identidad en el disco de perturbación.

La acción tiene términos $-q_e\log q_e+(d_e-1)(1-q_e)\log(1-q_e)$ y $-h_f\log h_f$.
Para el primero, $\log q_e=\log P_e+v_e$ sobre el soporte y
$\sum_e P_e|\log P_e|\le j\log j$ por la entropía de las filas.
Los demás términos cuestan en suma $O(\rho j)$ usando las cotas de cargas anteriores.
El factor local del predictor se calcula con $\log(1-d_e q_e)-d_e\log(1-q_e)$,
también de coste $O(\rho j)$. Se resta el valor basal y se conservan los signos del
predictor; las cotas son sobre normas. Con ello
\[
 \|\Lambda\|_\rho\le C\{j\log(ej)+\rho^2j\}.
\]
Para $\rho>1$ fijo inferior a dos se puede usar la cota con dos. Para $\rho$ admisible
variable el lado derecho es $\le C'j^2$. No exige una cota inferior por casilla positiva.
La identidad con los modelos de una línea respeta las incidencias unitarias: con $h=1-q$,
$h+T_d=(1-q)(1-(d-1)q)/(1-dq)$ y $T_d h/(h+T_d)=T_{d-1}$. Este pivote cancela
el cambio de factor local, y en la acción los términos unitarios también se cancelan.
